\% !TEX root = ../../buch.tex \% !TEX encoding = UTF-8 \%
\textbackslash section\{Die Taylor-Reihe als Grundlage\}
\textbackslash label\{pade:section:taylor\}
\textbackslash kopfrechts\{Taylor-Reihe\}

Die Padé-Approximation baut vollständig auf der Taylor-Reihe einer
Funktion auf. Ohne die Taylor-Reihe lässt sich keine Padé-Approximation
konstruieren, da sämtliche Koeffizienten der späteren rationalen
Funktion aus den Koeffizienten der Taylor-Reihe bestimmt werden. Bevor
das eigentliche Verfahren vorgestellt wird, sollen deshalb zunächst die
wichtigsten Eigenschaften der Taylor-Reihe erläutert werden.

Die Grundidee der Taylor-Reihe besteht darin, eine genügend oft
differenzierbare Funktion in der Umgebung eines bestimmten Punktes durch
eine Potenzreihe zu beschreiben. Dieser Punkt wird als Entwicklungspunkt
bezeichnet und üblicherweise mit \textbackslash(a\textbackslash)
notiert. Die allgemeine Taylor-Reihe einer Funktion
\textbackslash(f(x)\textbackslash) lautet
\textbackslash begin\{equation\} f(x) =
\textbackslash sum\_\{k=0\}\^{}\{\textbackslash infty\}
\textbackslash frac\{f\^{}\{(k)\}(a)\}\{k!\}(x-a)\^{}k.
\textbackslash label\{eq:taylor\_allgemein\}
\textbackslash end\{equation\}

Jeder Summand dieser Reihe enthält eine Ableitung der Funktion am
Entwicklungspunkt. Die erste Ableitung beschreibt die Steigung der
Funktion, höhere Ableitungen liefern Informationen über ihre Krümmung
und ihr weiteres Verhalten. Werden genügend viele Terme berücksichtigt,
kann das Verhalten einer Funktion in der Umgebung des
Entwicklungspunktes sehr genau beschrieben werden.

In vielen Anwendungen wird der Entwicklungspunkt auf
\textbackslash(a=0\textbackslash) gelegt. Man spricht dann von einer
\textbackslash emph\{Maclaurin-Reihe\}.
Gleichung\textasciitilde\textbackslash eqref\{eq:taylor\_allgemein\}
vereinfacht sich dadurch zu \textbackslash begin\{equation\} f(x) =
\textbackslash sum\_\{k=0\}\^{}\{\textbackslash infty\}
\textbackslash frac\{f\^{}\{(k)\}(0)\}\{k!\}x\^{}k.
\textbackslash label\{eq:maclaurin\} \textbackslash end\{equation\}

Ein besonders anschauliches Beispiel ist die Exponentialfunktion
\textbackslash{[} f(x)=e\^{}x. \textbackslash{]} Ihre Ableitung ist
wiederum die Exponentialfunktion selbst. Daraus folgt, dass sämtliche
Ableitungen an der Stelle \textbackslash(x=0\textbackslash) den Wert
\textbackslash(1\textbackslash) besitzen. Durch Einsetzen in
Gleichung\textasciitilde\textbackslash eqref\{eq:maclaurin\} erhält man
die bekannte Reihenentwicklung \textbackslash begin\{equation\} e\^{}x =
1 + x + \textbackslash frac\{x\^{}2\}\{2!\} +
\textbackslash frac\{x\^{}3\}\{3!\} +
\textbackslash frac\{x\^{}4\}\{4!\} + \textbackslash cdots.
\textbackslash label\{eq:exp\_reihe\} \textbackslash end\{equation\}

Da eine unendliche Reihe auf einem Computer nicht vollständig berechnet
werden kann, wird sie nach endlich vielen Summanden abgebrochen. Das
entstehende Polynom wird als Taylor-Polynom bezeichnet.

Für die Exponentialfunktion ergibt sich beispielsweise das
Taylor-Polynom zweiter Ordnung zu \textbackslash begin\{equation\}
T\_2(x) = 1 + x + \textbackslash frac\{x\^{}2\}\{2\}.
\textbackslash label\{eq:taylor2\} \textbackslash end\{equation\}

Mit steigender Ordnung werden mehr Terme der Taylor-Reihe
berücksichtigt, wodurch sich die Genauigkeit der Approximation im
Allgemeinen verbessert.

\textbackslash section\{Grenzen der Taylor-Reihe\}
\textbackslash label\{pade:section:grenzen\}
\textbackslash kopfrechts\{Grenzen der Taylor-Reihe\}

Obwohl Taylor-Polynome in vielen mathematischen Anwendungen sehr
erfolgreich eingesetzt werden, besitzen sie auch grundlegende
Einschränkungen. Sie approximieren eine Funktion nur in der
unmittelbaren Umgebung ihres Entwicklungspunktes besonders gut. Mit
zunehmendem Abstand nimmt der Approximationsfehler häufig deutlich zu.

Dieses Verhalten lässt sich am Beispiel der Exponentialfunktion
beobachten. In der Nähe von \textbackslash(x=0\textbackslash) stimmen
die Werte des Taylor-Polynoms nahezu mit der Exponentialfunktion
überein. Für grössere Werte von \textbackslash(x\textbackslash) wird
jedoch sichtbar, dass das Polynom das starke Wachstum der
Exponentialfunktion immer schlechter beschreibt.
Abbildung\textasciitilde\textbackslash ref\{fig:taylorvergleich\} zeigt
diesen Vergleich.

\%Bild 1, Teil1, (taylorvergleich), Taylorpolynom vs Exponentialfunktion

Die Ursache hierfür liegt im Aufbau eines Polynoms. Ein Polynom besteht
ausschliesslich aus Potenzen der Variablen und besitzt deshalb nur eine
begrenzte Flexibilität. Eigenschaften wie starkes exponentielles
Wachstum oder asymptotisches Verhalten lassen sich damit nur
eingeschränkt beschreiben.

Dies führt zur zentralen Fragestellung dieses Kapitels: Kann dieselbe
Information der Taylor-Reihe auf eine andere Weise genutzt werden, um
eine genauere Approximation zu erhalten? Genau an dieser Stelle setzt
die Padé-Approximation an. Anstatt die Koeffizienten der Taylor-Reihe
direkt zu einem Polynom zusammenzufassen, werden sie verwendet, um den
Zähler und den Nenner einer rationalen Funktion zu bestimmen. Dadurch
entsteht häufig eine deutlich genauere Näherung der ursprünglichen
Funktion.

Im nächsten Abschnitt wird dieses Verfahren vorgestellt und anhand der
Exponentialfunktion Schritt für Schritt hergeleitet.
